%-------------------------------------------------------------------------------
%   CHAPTER 1 - INTRODUCTION
%-------------------------------------------------------------------------------

\section{Удиртгал}
\label{sec:introduction}

\subsection{Асуудлын тулгамдсан байдал}

Монгол Улсын жижиг, дунд бизнесийн байгууллагууд (ЖДБ) хоорондоо бараа бүтээгдэхүүн захиалах, хүргэх үйл явцаа өдгөө хүртэл гол төлөв утас, мессеж, цаасан бичгээр гардаг~\cite{nso2023}. Захиалагч барааг өөрийн биеэр авч ирэх, эсхүл хүргэлтийг тусад нь зохион байгуулах шаардлага нь цаг алдагдуулж, зохион байгуулалтын алдаа үүсгэж, үйл явцыг хянахад төвөгтэй болгож байна.

Дижитал шилжилтийн нөхцөлд B2B харилцааг автоматжуулах эрэлт дэлхий нийтэд эрс нэмэгдэж~\cite{mckinsey2021delivery, grandview2023}, хүргэлтийн цахим платформууд ($\approx$4.5 тэрбум ам.доллар, 2023 он) гол хөдөлгөгч хүч болж байна~\cite{wang2016bigdata}. Монголын зах зээлд энэ чиглэлийн нэгдсэн шийдэл хараахан төлөвшөөгүй бөгөөд орон нутгийн тоглогчид (Papa Logistics гэх мэт) голдуу dispatcher-т тулгуурласан, нэг талтай загвартай хэвээр байна.

\subsection{Судалгааны асуулт}

Энэ ажил нь \textbf{operator-гүй, multi-tenant B2B хүргэлтийн платформыг одоогийн cloud BaaS технологи дээр найдвартай хэрхэн бүтээх вэ?} гэсэн ерөнхий асуултыг дэвшүүлж, гурван тодорхой дэд асуултыг шалгана:

\begin{itemize}
    \item \textbf{RQ1 — Concurrency:} PostgreSQL-ийн \texttt{FOR UPDATE SKIP LOCKED} механизм нь FCFS claim model-ыг race condition-гүй хэрэгжүүлж чадах уу?
    \item \textbf{RQ2 — Isolation:} Row Level Security (RLS) нь application-layer алдаанаас үл хамааран multi-tenant өгөгдлийн тусгаарлалтыг database түвшинд хангаж чадах уу?
    \item \textbf{RQ3 — Viability:} BaaS-mediated архитектур нь зөвхөн prototype биш, production-д хүрэх системийг хурдан хөгжүүлэх суурь болж чадах уу?
\end{itemize}

\subsection{Зорилго, зорилт}

Ажлын \textbf{зорилго} нь B2B хүргэлтийн operator-гүй, multi-tenant платформыг Supabase (PostgreSQL + PostgREST + RLS) технологи дээр боловсруулж, гурван дэд асуултыг туршилтаар хариулах явдал юм.

Дэвшүүлсэн \textbf{зорилтууд}:

\begin{enumerate}
    \item \textbf{Судалгаа, шинжилгээ:} Монгол Улс дахь B2B хүргэлтийн өнөөгийн байдлыг дүгнэх, Uber Freight, Papa Logistics зэрэг ижил төрлийн платформуудыг харьцуулах, архитектурын зөрүү ба design шаардлагыг тодорхойлох.
    \item \textbf{Архитектур, загварчлал:} Client layer (React Native/Expo + Next.js), BaaS layer (Supabase), data layer (PostgreSQL + RLS)-ээс тогтсон BaaS-mediated архитектур зохион бүтээх. 11 хүснэгт бүхий schema, ERD, sequence diagram, use case diagram боловсруулах.
    \item \textbf{Аюулгүй байдал:} RBAC, RLS, JWT authentication, bcrypt ePOD OTP-г нэгтгэсэн олон давхаргат security хэрэгжүүлэх.
    \item \textbf{Функционал модулиуд:} Захиалга, хүргэлтийн даалгавар, marketplace, realtime мэдэгдэл, ePOD, орлого, санхүүгийн тайланг бүрэн хэрэгжүүлэх.
\end{enumerate}

\subsection{Хамрах хүрээ ба хязгаарлалт}

Ажлын хамрах хүрээ нь \textbf{(i)} Монгол Улсын B2B хүргэлтийн orchestration, \textbf{(ii)} дунд зэргийн ачаалал (100 concurrent user хүртэл) бүхий staging орчинд хэрэгжсэн техникийн нотолгоо, \textbf{(iii)} single-region Supabase Pro tier тохиргоогоор хязгаарлагдана. Бодит GPS location tracking, төлбөрийн бодит integration (QPay, SocialPay), ML recommendation нь энэ ажлын хүрээнд биш — цаашдын хөгжүүлэлтийн чиглэл юм (§\ref{subsec:results_summary}).

\subsection{Ажлын бүтэц}

Үлдэж буй хэсэг нь дараах бүтэцтэй. \textbf{Бүлэг 2}-т last-mile логистик, two-sided marketplace, multi-tenant архитектур, concurrency control, BaaS-ын онолын суурь болон Uber Freight, Papa Logistics-тэй харьцуулсан судалгааг танилцуулна. \textbf{Бүлэг 3}-т системийн архитектур, өгөгдлийн загвар, аюулгүй байдлын бодлого, функционал шаардлагыг дэлгэрэнгүй тайлбарлана. \textbf{Бүлэг 4}-т гүйцэтгэсэн туршилтын үр дүн, харьцуулалт болон дүгнэлтийг оруулсан.
